\documentclass[11pt]{article}
\usepackage[margin=1in]{geometry}
\usepackage[T1]{fontenc}
\usepackage[utf8]{inputenc}
\usepackage{lmodern,microtype,xspace}
\usepackage{amsmath,amssymb,mathtools,booktabs,enumitem,hyperref}
\newcommand{\HsH}{H(s)H\xspace}
\newcommand{\mcite}[1]{\textsuperscript{\tiny\textsf{[#1]}}}
\newcommand{\Hrel}{\mathsf{H}}
\newcommand{\Forest}{\mathcal F}
\title{\textbf{The Forest: Physical Information Has a Path}\\[0.3em]\large Worldtube Locality, Hagalaz Relations, and Threefold Closure in SAT/\HsH}
\author{Nathan McKnight\\\small Working draft with LLM-assisted formalization}
\date{21 September 2026}

\begin{document}
\maketitle

\begin{abstract}
We develop a conditional bridge among four ideas: (i) information \emph{about an actual thing} requires a realized physical ancestry connecting that thing to a record; (ii) what appears nonlocal in a three-dimensional projection may be local with respect to native four-dimensional worldtube connectivity; (iii) Hagalaz supplies a candidate relation language for displacement, scale, frame rotation, and internal state; and (iv) a threefold circle/helix construction places the historical SAT quantity $B=3/(4\pi)$ directly beside the $120^\circ$ tangent hierarchy. The first claim is formulated using mutual information and data processing. SAT/\HsH-specific mechanisms remain conditional model hypotheses, not established physics.
\end{abstract}

\paragraph{Status.}
``Physical'' is used operationally for a realized, difference-bearing process or state in the world: observation, photons, ink, voltages, neural activity, stored bits, sound, and any other actual carrier. Arbitrary encoding changes the mapping between carrier and meaning; it does not make the carrier or transmission nonphysical. No metaphysical claim about the ultimate substrate is required.

\section{Physical information has a path}

\textbf{Physical Information Principle.} If a record contains information about a particular actual referent, there exists a realized physical dependency path connecting that referent, or a physical trace generated by it, to the record.\mcite{P:N21a}

Write a realized ancestry schematically as
\begin{equation}
X\longrightarrow M_1\longrightarrow\cdots\longrightarrow M_k\longrightarrow R,
\end{equation}
where $X$ is a contingent referent state, $M_i$ are physical mediators/encodings/measurements/memories, and $R$ is a record. For a Markov segment, the data-processing inequality gives
\begin{equation}
I(X;M_{i+1})\le I(X;M_i).
\end{equation}
If mathematics produces $Y=f(R,K)$ from a record $R$ and prior knowledge $K$, then
\begin{equation}
I(X;Y)\le I(X;R,K).
\end{equation}
Mathematics can reorganize, expose, sharpen, or discard information present in its inputs; it cannot manufacture referential information about $X$ when the physically instantiated inputs contain none. The empirical content of $K$ has its own physical ancestry.

A useful operational form is the \emph{alter-the-thing test}: change the target property while holding fixed every physically available state on which the alleged knower actually depends. If the record cannot change, the record does not carry information about that changed property.

\section{The Forest and the meaning of locality}

For a record $R$, let its information-ancestry tree be the physically realized dependency subgraph supporting its referential content. The union of such overlapping structures is the \emph{Forest},
\begin{equation}
\Forest=(V,E),
\end{equation}
where vertices are realized states/events/extended sections and edges are allowed difference-bearing relations or channels.

This turns locality into a property of the theory's primitive adjacency relation, not merely a synonym for small projected spatial separation.

SAT makes a stronger wager: apparent nonlocality may be physically transmitted in a direction transverse to ordinary time propagation along an extended worldtube.\mcite{P:N21b} Let $\sim_{3+1}$ denote ordinary projected spacetime locality and $\sim_W$ native worldtube adjacency. The proposal is
\begin{equation}
A\not\sim_{3+1}B\qquad\text{while possibly}\qquad A\sim_W B.
\end{equation}
Thus apparent nonlocality is not pathlessness; it is a mismatch between projected locality and native physical connectivity.

This does not evade Bell's theorem by vocabulary. A Bell-local hidden-variable model factorizes as
\begin{equation}
P(a,b\mid x,y,\lambda)=P(a\mid x,\lambda)P(b\mid y,\lambda).
\end{equation}
If SAT is to reproduce Bell correlations using a real worldtube-transverse channel, the channel must change the relevant causal/factorization structure relative to ordinary $3+1$ locality while remaining compatible with operational no-signalling.

\section{Hagalaz as an edge language}

A current Hagalaz reconstruction represents a framed object by
\begin{equation}
\mathcal S_i=(c_i,D_i,F_i;\mathcal P_i),
\end{equation}
with center $c_i$, positive scale $D_i$, orthonormal frame $F_i$, and internal state $\mathcal P_i$. Its narrow pairwise relation is
\begin{equation}
\Hrel_{ij}=(\Delta c_{ij},\sigma_{ij},Q_{ij}),\quad
\Delta c_{ij}=c_j-c_i,\quad
\sigma_{ij}=\frac{D_j}{D_i},\quad
Q_{ij}=F_jF_i^{-1}.
\end{equation}
The elementary relation identities have been independently checked, but this is not yet a Hagalaz dynamical law.\mcite{P:HAG-R}\mcite{P:HAG-T}

For a closed three-node cycle,
\begin{align}
\Delta c_{12}+\Delta c_{23}+\Delta c_{31}&=0,\\
\sigma_{12}\sigma_{23}\sigma_{31}&=1,\\
Q_{31}Q_{23}Q_{12}&=I.
\end{align}
For a longer path $\gamma$, define rotational holonomy
\begin{equation}
Q_\gamma=Q_{N1}Q_{N-1,N}\cdots Q_{12}.
\end{equation}
This is a natural local-to-global object: repeated local frame relations can become a global closure condition.

\section{Three circles and the candidate appearance of $B$}

Arrange $n$ equal tangent circles of radius $r_f$ cyclically around a central aperture. Their centers form a regular $n$-gon and the inscribed aperture radius is
\begin{equation}
h_n=r_f\left[\csc\left(\frac{\pi}{n}\right)-1\right].
\end{equation}
For $n=2$, $h_2=0$: the aperture degenerates to the original two-circle kissing point. At the central kissing points, consistently oriented tangent directions advance by
\begin{equation}
\Delta_n=\frac{2\pi}{n}.
\end{equation}
Hence the ladder begins
\begin{equation}
\Delta_2=\pi,\qquad \Delta_3=\frac{2\pi}{3},\qquad \Delta_4=\frac{\pi}{2},\ldots
\end{equation}
The historical SAT quantity
\begin{equation}
B=\frac{3}{4\pi}
\end{equation}
is well attested in the archive.\mcite{P:B-HIST}

If one defines a two-sided tangent normalization
\begin{equation}
B_n^{(\mathrm{tan})}\equiv\frac{1}{2\Delta_n},
\end{equation}
then exactly
\begin{equation}
B_n^{(\mathrm{tan})}=\frac{n}{4\pi},\qquad
B_3^{(\mathrm{tan})}=\frac{3}{4\pi}=B.
\end{equation}
The equality is exact; the interpretation is not. In particular, the factor of two must be geometrically derived rather than selected because it produces $B$. Candidate sources include the two orientations of a tangent line, two chirality branches, or another paired tangent construction. Until that factor is forced, this is a candidate identification, not a derivation of the historical role of $B$.

\section{Triswale, helix pitch, and chirality}

Take the three circle planes as equatorial planes with coplanar origins and twist axes separated by $120^\circ$. Cooperative signed twist produces a threefold central saddle (the \emph{triswale}). Near the center, the lowest smooth threefold mode is
\begin{equation}
z(\rho,\phi)=A_3\rho^3\cos\bigl(3(\phi-\phi_0)\bigr)+\mathcal O(\rho^4).
\end{equation}
The signed cubic coefficient $A_3$ is a natural candidate order parameter because the canonical cubic saddle has no quadratic curvature at the center.

Let $\theta_4$ be the helix tangent angle measured from the axial/common-normal direction. For helix radius $a$ and signed pitch per turn $p$,
\begin{equation}
\tan\theta_4=\frac{2\pi a}{p},\qquad p=2\pi a\cot\theta_4.
\end{equation}
This reference convention must remain explicit because $\theta_4$ has undergone semantic drift in SAT/\HsH.\mcite{P:T4-DRIFT}

Define
\begin{equation}
s\equiv\frac{p}{2\pi a}=\cot\theta_4.
\end{equation}
In the ideal zero-thickness limit, $s>0$ and $s<0$ label opposite helix chiralities, while $s=0$ is the zero-pitch densely wound limit. If triswale response is smooth and mirror-odd,
\begin{equation}
A_3(-s)=-A_3(s),\qquad A_3(s)=\kappa_3s+\mathcal O(s^3),
\end{equation}
which is a testable ansatz, not yet an exact contact solution.

Finite filament thickness must be handled separately: neighboring turns can contact before centerline pitch reaches zero, and contact, fusion, and topological reconnection are not the same event. The finite-thickness variable should therefore be true minimum surface-to-surface rung clearance, derived from tube geometry.

\section{Local twist to global closure}

The candidate throughline is now
\begin{equation}
\boxed{\begin{gathered}
\text{threefold packing}\to\text{tangent invariant}\to B\to\text{signed twist}\\
\to Q_{ij}\to Q_\gamma\to\text{global chirality/closure}.
\end{gathered}}
\end{equation}
Every arrow is to be independently derived or rejected.

If the large-scale geometry is toroidal, repeated local relations may close only in discrete winding sectors. This is where ``number of twists in the donut'' becomes a mathematical question: define the local step, path, composition law, and closure condition. A further rule is still required before any particle census can be inferred:
\begin{equation}
N_{\mathrm{particles}}\stackrel{?}{=}f(N_{\mathrm{winding}}).
\end{equation}
That mapping is open.

The information principle re-enters here. If SAT nonlocality is a real transverse worldtube transmission, the same native geometry that determines $Q_{ij}$ and $Q_\gamma$ should determine which Forest vertices are physically adjacent for information transfer. Information and geometry are not separate ontologies: an information path is a path in the physical geometry.

\section{Three-state proxies and asymmetry questions}

A three-circle model admits a provisional binary label on each strand. If $u$ and $d$ are used only as opposite local twist labels,
\begin{equation}
uud\sim(+,+,-),\qquad udd\sim(+,-,-),
\end{equation}
with net signed local chirality
\begin{equation}
q_\chi(uud)=+1,\qquad q_\chi(udd)=-1.
\end{equation}
A globally achiral geometry makes these mirror sectors symmetric at this level. A global toroidal chirality $\chi_U$ allows an odd coupling such as
\begin{equation}
E_\chi=\lambda\chi_Uq_\chi+\cdots.
\end{equation}
This has the right symmetry to split the proxy sectors, but it is not a proton--neutron mass derivation and does not identify the labels with Standard Model quarks.

Likewise, matter--antimatter asymmetry requires more than a signed geometric preference. Geometry may supply an order parameter or bias; a physical baryogenesis mechanism would still require the appropriate number-changing, symmetry-violating, and non-equilibrium dynamics.

\section{Immediate tests}

\begin{enumerate}[leftmargin=1.6em]
\item Define SAT/\HsH native worldtube adjacency and the physical degree(s) of freedom carrying time-perpendicular influence.
\item Give Hagalaz a transport/dynamical law rather than stopping at relation coordinates.
\item Solve the exact three-circle/finite-tube cooperative twist and derive $A_3(s)$.
\item Derive or reject the factor of two in $B=1/(2\Delta_3)$ by geometry and historical genealogy.
\item Compose the local relation around the torus and compute the actual holonomy/winding closure law.
\item Only then test $uud/udd$ proxies, proton/neutron-scale ratios, matter/antimatter bias, or particle-count mappings.
\item Construct an explicit Bell/CHSH probability model for the worldtube-transverse channel and determine exactly which ordinary Bell-locality assumption is altered while preserving operational no-signalling.
\end{enumerate}

\section{Conclusion}

The throughline is not that information is abstract and geometry is physical. Information about an actual thing is itself a physical relation: a realized ancestry of observation, recording, transformation, and transmission. Mathematics can operate on that ancestry but cannot conjure the referent into its symbols.

SAT/\HsH then makes an additional wager: ordinary spatial separation may fail to capture native physical adjacency because extended worldtubes possess a transverse connectivity along which real influence can propagate. If that wager is correct, nonlocality is not pathlessness; it is locality in a geometry whose primitive edges differ from those of the ordinary projection.

Hagalaz is a candidate language for those edges. Threefold circle geometry supplies a $120^\circ$ local structure and an exact candidate appearance of $B=3/(4\pi)$. Signed triswale/helix twist supplies a local chirality variable. Composition around a toroidal path supplies the natural local-to-global object: holonomy.

Thus the present Forest hypothesis is
\begin{equation}
\boxed{\begin{gathered}
\text{information path}\equiv\text{physical path}\subset\text{worldtube relation geometry};\\
\text{local geometry composes into global closure}.
\end{gathered}}
\end{equation}
Everything after the semicolon remains to be earned by explicit geometry and dynamics.

\section*{Project provenance resolver}
\begin{description}[leftmargin=1.7cm,style=nextline]\small
\item[P:N21a] Nathan-direct correction, current 2026-09-21 conversation: information about an actual thing requires a physical observation/transmission ancestry; arbitrary encoding does not make its realization nonphysical. Raw conversation/message ID pending archive backfill.
\item[P:N21b] Nathan-direct statement, same conversation: SAT foundational assumption that apparent nonlocality is physically transmitted in a time-perpendicular direction along worldtubes; if wrong without a close substitute, SAT is likely a nonstarter.
\item[P:HAG-R] `WORKSPACES/MERIDIAN/HAGALAZ_FRAMED_HYPERSPHERE_NOTE_2026-09-20.md`.
\item[P:HAG-T] `WORKSPACES/MERIDIAN/HAGALAZ_REPRESENTATION_TEST_2026-09-20.md`.
\item[P:B-HIST] `SAT_THEORY_ARCHIVE_2023-25/[[SAT26 TOOLBOX]]/SAT26 MATH ROUNDUP.txt`, used only as historical provenance for $B=3/(4\pi)$.
\item[P:T4-DRIFT] `WORKSPACES/MERIDIAN/HANDOFF_NWTF_THETA4_SEMANTIC_DRIFT_2026-09-20.md`.
\end{description}

\begin{thebibliography}{9}
\bibitem{CoverThomas} T. M. Cover and J. A. Thomas, \emph{Elements of Information Theory}, 2nd ed. (Wiley, 2006), Ch. 2. DOI: 10.1002/047174882X.
\bibitem{Landauer} R. Landauer, ``Information is physical,'' \emph{Physics Today} \textbf{44}(5), 23--29 (1991). DOI: 10.1063/1.881299.
\bibitem{Bell} J. S. Bell, ``On the Einstein Podolsky Rosen paradox,'' \emph{Physics Physique Fizika} \textbf{1}, 195--200 (1964). DOI: 10.1103/PhysicsPhysiqueFizika.1.195.
\bibitem{CHSH} J. F. Clauser, M. A. Horne, A. Shimony, and R. A. Holt, ``Proposed experiment to test local hidden-variable theories,'' \emph{Physical Review Letters} \textbf{23}, 880--884 (1969). DOI: 10.1103/PhysRevLett.23.880.
\end{thebibliography}
\end{document}
